\begin{lemma}[The $A$-bound]
Let $a,b,s,q\in\mathbb{R}$ satisfy
\[
  0<a,\qquad a\le s,\qquad q=a^2+b^2,\qquad s^2+b^2=1 .
\]
Then
\[
  0\ \le\ \frac{b^2(s-a)}{q}\ \le\ \frac{(1-a)^2}{2a}.
\]
\end{lemma}
\begin{proof}
\emph{Step 0: $q>0$ and $s+a>0$.} Since $a>0$ we have $a^2>0$, and $b^2\ge0$, so
$q=a^2+b^2\ge a^2>0$. From $0<a\le s$ we get $s+a\ge2a>0$.

\emph{Step 1: the lower bound.} We have $b^2\ge0$, $s-a\ge0$ by hypothesis, and
$q>0$ by Step 0, so $b^2(s-a)\ge0$ and hence $b^2(s-a)/q\ge0$.

\emph{Step 2: two identities.} First,
\begin{equation}\label{eq:ab-sa}
  s^2-a^2=(s^2+b^2)-(a^2+b^2)=1-q ,
\end{equation}
using both $s^2+b^2=1$ and $q=a^2+b^2$. Second, $b^2=q-a^2$, again from
$q=a^2+b^2$. Therefore, multiplying numerator and denominator by $s+a>0$ and
using \eqref{eq:ab-sa},
\begin{equation}\label{eq:ab-rewrite}
  \frac{b^2(s-a)}{q}
  =\frac{b^2(s-a)(s+a)}{q\,(s+a)}
  =\frac{b^2(s^2-a^2)}{q\,(s+a)}
  =\frac{(q-a^2)(1-q)}{q\,(s+a)}
  =\frac{1+a^2-q-a^2/q}{s+a},
\end{equation}
the last equality because multiplying $1+a^2-q-a^2/q$ by $q>0$ gives
\[
  q+a^2q-q^2-a^2=q(1-q)+a^2(q-1)=(1-q)(q-a^2),
\]
so that $1+a^2-q-a^2/q=(1-q)(q-a^2)/q$. This is the paper's rewriting of
$b^2(s-a)/q$.

\emph{Step 3: the numerator is at most $(1-a)^2$.} Since $q>0$,
\[
  q+\frac{a^2}{q}-2a=\frac{q^2-2aq+a^2}{q}=\frac{(q-a)^2}{q}\ \ge\ 0,
\]
i.e. $q+a^2/q\ge2a$. Hence
\[
  1+a^2-q-\frac{a^2}{q}
  =1+a^2-\Bigl(q+\frac{a^2}{q}\Bigr)
  \ \le\ 1+a^2-2a=(1-a)^2 .
\]

\emph{Step 4: the upper bound.} Dividing the inequality of Step 3 by the
positive number $s+a$ (Step 0) and using \eqref{eq:ab-rewrite},
\[
  \frac{b^2(s-a)}{q}=\frac{1+a^2-q-a^2/q}{s+a}\ \le\ \frac{(1-a)^2}{s+a}.
\]
Finally $(1-a)^2\ge0$ and $0<2a\le s+a$ by Step 0, so
$(1-a)^2/(s+a)\le(1-a)^2/(2a)$. Combining the last two displays gives
$b^2(s-a)/q\le(1-a)^2/(2a)$, which together with Step 1 is the assertion.
\end{proof}
